%!TEX root = ../template.tex
%% ====================================================================
%% Capitulo 5 -- Resultados I --- A plataforma
%% Primeira versao completa 2026-09-08. Linguagem descritiva.
%% Fontes da vistoria (2026-09-08): ~/tese-extraccao (modelo v22 de
%% 2026-07-31: modelo.bim, medidas, queries, particoes, metadados, 25 TSV,
%% .pbix com 4 paginas, 33 prints), notas do Mind Palace de 2026-07-01/03
%% (mapeamento e validacao a 3 fontes) e memoria do projeto (2026-07-05).
%% Marcas %% CONFIRMAR assinalam factos a confirmar pelo autor.
%% Figuras: prints em ~/tese-extraccao/prints_dashboard (recortar e
%% anonimizar antes de incluir); sugeridas 4, uma por pagina.
%% ====================================================================

\chapter{Resultados I --- A plataforma}
\label{cha:plataforma}

O primeiro resultado desta dissertação é um instrumento. Este capítulo
descreve o que foi entregue --- a arquitetura da plataforma de
monitorização energética que está em uso na refinaria ---, prova que a
camada de dados diz a verdade, apresenta o \emph{dashboard} que a torna
legível, e declara o que ficou de fora. Tudo o que o
Capítulo~\ref{cha:diagnostico} conclui sobre as \emph{baselines} é uma
consulta ao modelo de dados aqui descrito; se este capítulo falhasse, o
seguinte não teria base. A ordem do capítulo é a ordem da confiança: só
depois de se ter mostrado que os números batem certo é que faz sentido
mostrar como se leem.

\section{Arquitetura entregue}
\label{sec:arquitetura}
% [REV-2026-09-13][C5-01][P2][ESTRUTURA]
\revisaotese{C5-01}{P2}{ESTRUTURA}
{Separar resultado entregue de características da implementação}
{O inventário da arquitetura repete o Cap. 4. Aqui bastam funcionalidades,
âmbito coberto e evidência de entrega/uso: versão, data, modo de publicação,
cadência de atualização e responsável pela manutenção. Tempo de refresh,
latência dos dados e desempenho em tarefas concretas são evidências
melhores de utilidade que número de medidas ou tamanho do ficheiro.
Se não foram medidos, limitar a conclusão a implementação funcional.}%
% [/REV-2026-09-13]

\paragraph{O que é.} A plataforma é um modelo tabular em Power BI ---
\emph{Energy Performance Monitor}, versão 22, de 31 de julho de 2026 ---
com um esquema em estrela de 32 tabelas, 243 colunas, 13 relações, 159
medidas e 21 consultas de transformação, que cobre 16 unidades processuais
em três fábricas e nos extras (enxofres e \emph{flares}), sete vetores
energéticos, resolução diária, de 1 de janeiro de 2020 ao dia anterior a
cada atualização. O modelo ocupa 68,7~MB e importa os dados; o historiador
é consultado através da plataforma de dados da refinaria, os balanços e os
relatórios de eletricidade lidos das pastas de rede onde o Processamento de
Dados os mantém, e os planos de fabrico do arquivo de planeamento. As
consultas ao historiador são geradas a partir das tabelas de mapeamento
no momento da atualização; nenhuma \emph{tag} está escrita no código.

\paragraph{Três decisões que o definem.} A primeira é a que a
Secção~\ref{sec:modelo-dimensional} descreve: os três canais de origem
coexistem na tabela de factos e uma regra de precedência declarada escolhe
por linha, com a origem registada --- é o que torna a validação da secção
seguinte possível, porque um dia com balanço e \emph{tag} em simultâneo é
um dia com duas leituras da mesma grandeza, processadas por caminhos
diferentes, que podem ainda assim partilhar o medidor. A segunda é que a
\emph{baseline} é materializada no modelo: a tabela de coeficientes por
unidade, vetor e ano é calculada uma vez por atualização e as medidas de
resíduo, banda e excesso leem-na por pesquisa, em vez de reestimarem a
regressão em cada visual --- o que tornou o \emph{dashboard} responsivo e,
mais importante, tornou a \emph{baseline} exportável e auditável como uma
tabela. A terceira é que todos os parâmetros que um utilizador pode querer
mudar --- ano de referência da \emph{baseline}, métrica (toneladas,
percentagem ou desvios-padrão), janela (último dia, 7 dias, 30 dias,
acumulado), unidade, vetor, período --- são seletores do modelo e não
cópias de páginas: as 159 medidas são genéricas em relação ao contexto de
filtro, e a mesma medida serve as quatro páginas.

\paragraph{Como está organizado.} As medidas vivem em tabelas de domínio
--- consumos e conversões gerais, \emph{baselines}, resíduos e bandas,
séries temporais e acumulados, índice de intensidade, seletores e cartões
--- para que quem mantém o modelo encontre cada cálculo pelo que ele faz.
Os visuais que o Power BI nativo não oferece (o diagrama de dispersão com
banda e cor temporal, o mapa de calor de calendário, a cascata, a matriz
de prioridades, os cartões de indicador) são declarados em Vega-Lite ou em
HTML dentro do próprio relatório, sem dependências externas, e leem as
mesmas medidas. Seis marcadores guardam as vistas de arranque de cada
página.

\paragraph{Evolução.} A primeira arquitetura, de março de 2026, ligava-se
ao historiador por consulta direta, usava um limiar de carga único e
estimava a regressão em tempo de consulta; a versão de julho, aqui
descrita, importa, lê o limiar de cada unidade da sua dimensão,
materializa a \emph{baseline} e reconcilia três canais. Entre as duas
versões o mapeamento mestre chegou à sua vigésima segunda revisão, e a
validação da secção seguinte motivou a maior parte das últimas.
%% CONFIRMAR: publicacao (workspace Power BI Service? cadencia de atualizacao?
%% utilizadores — equipa TEC de monitorizacao) — a memoria do projeto diz
%% "em producao"; precisar onde e desde quando.

\section{Validação da camada de dados}
\label{sec:validacao-dados}

Um \emph{dashboard} que apresenta números errados com confiança é pior do
que nenhum \emph{dashboard}. Antes de qualquer visual ser dado por
terminado, a camada de dados foi validada a três fontes: para cada
unidade, vetor e mês de janeiro a maio de 2026, o valor do
\emph{dashboard} foi comparado com o relatório oficial de consumo emitido
por cada fábrica (o ``Envio'') e com uma reconstrução independente a
partir do mapeamento mestre aplicado diretamente às folhas de medidores do
balanço. A amostra são 400 células: as combinações unidade--vetor do
relatório nos cinco meses, incluindo as que não têm consumo, contadas à
parte como vazias. Uma célula coincide quando a diferença absoluta entre o
\emph{dashboard} e as outras duas leituras não excede 2~t ou 0,2\% do valor,
o que for maior; o sinal da diferença regista-se para a análise de
consistência descrita adiante. As três leituras não são três medições
independentes. O relatório e a reconstrução partem do mesmo balanço e dos
mesmos medidores, e o \emph{dashboard} lê esse balanço por outro código:
a independência é do processamento, não da fonte. Uma divergência pode por
isso vir de uma fronteira ou de uma definição diferentes, e não de um valor
fisicamente errado, e uma coincidência confirma a contabilização, não a
medição física. O método existe para encontrar a causa de cada
divergência.

A Figura~\ref{fig:validacao-mapeamento} acompanha essa progressão versão a versão.

\begin{figure}[tbp]
  \centering
  \includegraphics[width=0.86\textwidth]{gerado/cap5-validacao-progressao}
  \caption[Progressão da validação do mapeamento]%
  {Progressão da validação do mapeamento, versão a versão. Cada barra reparte as mesmas 400 células unidade--vetor--mês da amostra de validação em coincidentes, divergentes e sem consumo, para que os vazios não sejam lidos como sucessos. As quatro primeiras barras são estados verificados célula a célula; a barra da v15.1, tracejada, é o resultado \emph{esperado} do fecho das últimas cinco divergências, não uma verificação observada. As causas de cada divergência e o número de casos que explicam estão descritas no texto que se segue.}
  \label{fig:validacao-mapeamento}
\end{figure}
\paragraph{O ponto de partida.} Na versão 13 do mapeamento, 242 células
coincidiam, 64 estavam vazias (combinações unidade--vetor sem consumo) e
94 divergiam. O trabalho de julho de 2026 foi explicar as 94, uma a uma,
até nenhuma ficar classificada como inconsistente. Ficaram explicadas por
sete causas, e o resultado mais instrutivo da validação é onde essas
causas estavam: em todas as três fontes.

\paragraph{Erros no motor.} O canal de balanço não tinha mecanismo de
inclusão: a coluna de factor do mapeamento aplicava-se só ao canal de
\emph{tags} (conversão de kg/h para t/h), e as dezoito linhas com factor
zero --- circuitos internos, medidores partilhados --- entravam por
inteiro, enquanto uma linha de crédito (factor $-1$) não era subtraída.
Duas unidades apareciam com o dobro do vapor real (48\,461 contra
26\,019~t num nível; 160\,322 contra 80\,696~t noutro, em cinco meses); a
correção, uma coluna nova de factor de inclusão para o canal de balanço,
moveu 251~kt de vapor em módulo sem tocar no canal de \emph{tags}. Um
filtro de nomes de folhas que só aceitava as fábricas excluía as folhas
dos enxofres e dos \emph{flares}: essas unidades nunca entravam na tabela,
e o mapeamento --- que já estava certo --- não podia ser validado. Um
ficheiro de segurança deixado na pasta de leitura dos relatórios de
eletricidade era lido ao vivo e somado ao definitivo, duplicando quatro
unidades em janeiro. E, antes destes, o desvio de colunas do balanço de
utilidades (Secção~\ref{sec:fontes}) sub-reportava os três níveis de vapor
em 5 a 7\% em todas as unidades --- o único destes erros que era
universal e silencioso.

\paragraph{Erros no mapeamento.} Uma referência de medidor em falta na
produção de vapor de 3~bar da destilação atmosférica produzia um défice
todos os meses --- 29,6; 61,1; 33,3; 9,4; 0,8~t --- que passava por baixo
da tolerância relativa em quatro dos cinco meses; foi encontrada por
procura do perfil mensal do défice entre as 481 referências da folha de
medidores, com coincidência a 0,1~t nos cinco meses. A alocação da
produção de vapor entre as duas unidades de vácuo seguia uma convenção
diferente da do controlo processual; das oito fórmulas oficiais, seis
foram confirmadas a $\pm0{,}5$~t e duas rejeitadas e devolvidas ao
controlo processual com a evidência (uma delas referia uma linha de
outra versão do ficheiro). Duas referências de fuel gás dos enxofres
tinham grafias que nunca existiram na folha --- um hífen no sítio errado,
uma letra a mais.

\paragraph{Erros na fonte oficial.} No relatório cumulativo de uma
unidade, um bloco de linhas de fevereiro tinha a coluna colada de maio:
produção de vapor a 24~bar 3\,264,0 no relatório contra 2\,215,4 nas
medições; consumo de vapor a 3~bar 1\,120,0 contra 2\,430,7; energia
elétrica 828,9 contra 569,1~MWh. O consumo específico do mês parecia
validar o valor colado --- porque era recalculado sobre ele com a carga
verdadeira, que escapara à colagem; o que o denunciou foi a igualdade ao
decimal entre a célula de fevereiro e a de maio. O Processamento de Dados
confirmou e corrigiu o relatório; as três células fecharam a três fontes
ao decimal. É o episódio que inverte a hierarquia habitual: a
``verdade-terreno'' estava errada e o \emph{dashboard} certo, e só uma
validação com três leituras processadas por caminhos diferentes o podia
mostrar.

\paragraph{O fecho.} Contando à parte as células vazias, a progressão foi
de 242 células numéricas coincidentes na versão 13 (com 64 vazias) para 263
depois do factor de inclusão e 327 com a versão 15 (com 65 vazias; uma
combinação deixou de ter consumo entre as duas versões), 330 com o
relatório de fevereiro corrigido e, com a versão 15.1, a correção das
últimas cinco, nas referências de fuel gás dos enxofres. As 94 divergências
iniciais ficaram todas explicadas. O registo de 3 de julho de 2026 dá o
fecho de 400 de 400 como o esperado na atualização seguinte com a
versão~15.1; essa execução não ficou registada com data, e o fecho completo
é, nesta versão do texto, um resultado esperado e não observado. A
validação cobre janeiro a maio de 2026 e não se estende por si aos anos de
2020 a 2025 que o Capítulo~\ref{cha:diagnostico} avalia
(Secção~\ref{sec:modelo-dimensional}).
%% CONFIRMAR: o 400/400 foi observado no refresh seguinte a v15.1 (a nota
%% de 2026-07-03 diz "esperado"); se sim, indicar a data.
O fuel gás foi validado a duas fontes, por decisão de âmbito, e a
ressalva fica escrita: o \emph{dashboard} coincide com o relatório porque
ambos leem o mesmo balanço de produção; é consistência entre fontes
dependentes, não verificação, e se o balanço vier errado erram os dois
sem aviso.

% [REV-2026-09-17][C5-06][P1][EVIDENCIA]
\revisaotese{C5-06}{P1}{EVIDENCIA}
{Confirmar se o 400/400 da versão 15.1 foi observado}
{O parágrafo anterior deixou de afirmar o 400/400 como resultado observado:
a nota de 2026-07-03 diz ``esperado no próximo refresh'' e não há registo
da execução. Se a validação foi corrida depois da 15.1 (ou sobre a v22 de
2026-07-31), indicar data e ficheiro e repor a afirmação aqui, no
Cap.~8 (Construir) e na Tabela~8.1. Confirmar também a unidade da tolerância
de 2 na eletricidade (MWh?).}%
% [/REV-2026-09-17]
\paragraph{A carga.} A carga foi validada por outro caminho: a soma dos
componentes de carga de cada unidade contra a linha de carga do balanço,
em 118 de 120 células (24 unidades, cinco meses) a $\pm0{,}1$~t; as duas
que não fecham são definição, não erro --- a linha de carga da destilação
atmosférica exclui o resíduo atmosférico reprocessado, que em fevereiro
valeu 9\% da carga do mês. Ficou identificada uma questão que o
\emph{dashboard} não resolve sozinho: os relatórios oficiais e o balanço
usam duas definições de carga (hidrocarbonetos ao nível do componente,
excluindo hidrogénio, limpezas, \emph{quench} e retornos, contra a linha
total do balanço), ambas reproduzíveis, e a escolha muda os consumos
específicos --- 2,4\% numa unidade só pelo hidrogénio, um factor de sete
noutra. É uma decisão de definição de indicador, registada como pendente
e devolvida a quem define o indicador.

A validação usa as 24 unidades do balanço; o \emph{dashboard} agrega-as em 16,
e a destilação atmosférica do \emph{dashboard} junta à coluna as secções
auxiliares declaradas na sua composição. A carga que entra nos indicadores
e na regressão dessa unidade é a linha de carga do balanço, que exclui o
resíduo atmosférico reprocessado e os \emph{slops}. A escolha pesa pouco na
avaliação que se segue, mas não é nula: com a carga total da coluna, 15 dos
2\,192 dias de 2020 a 2025 passariam o limiar, e o declive anual do fuel gás
mudaria entre $-2{,}5\%$ e $+3{,}4\%$; nos vapores, cujos declives são
próximos de zero, as mudanças relativas são grandes e as absolutas
pequenas. O Capítulo~\ref{cha:diagnostico} replica a convenção em produção;
replicá-la não diz se é a fisicamente correta.

\paragraph{O que a validação ensinou.} Quatro lições passaram a regras
da plataforma. Uma tolerância relativa em números grandes esconde termos
pequenos e sistemáticos: o segundo teste obrigatório é a consistência do
resíduo --- se a diferença tem sempre o mesmo sinal e nunca é nula, é
estrutural, por mais pequena que seja. Um filtro de ficheiros por prefixo
é uma armadilha (``fábrica I'' é prefixo de ``fábrica II''); os filtros
passaram a exigir o nome canónico completo e a recusar duplicados em vez
de os somar. Os ficheiros de segurança nunca ficam na pasta de leitura. E
a validação a três fontes não é um passo de projeto: é um teste de
aceitação permanente, que a plataforma repete a cada fecho mensal, porque
apanhou erros nas três fontes e nenhuma delas, isolada, era fiável.

\section{O \emph{dashboard}}
\label{sec:dashboard}
% [REV-2026-09-13][F09][P2][FIGURA]
\revisaotese{F09}{P2}{FIGURA}
{Dois screenshots anotados no corpo, quatro páginas no manual}
{Usar (1) vista da refinaria com percurso unidade -\textgreater{} vetor e (2) detalhe
de baseline com reta, dados, banda e acumulado; colocar cada imagem
após o respetivo parágrafo. Recortar ao canvas e anotar 3--5 elementos
referidos no texto; incluir data, filtros, janela e referência visíveis.
Usar capturas reais verificadas de prints\_dashboard, sem inventar estados.
Mover capturas integrais das quatro páginas e instruções para o apêndice do manual do dashboard.
Uma figura por página só se os rótulos ficarem legíveis à dimensão final.}%
% [/REV-2026-09-13]

O \emph{dashboard} tem quatro páginas, desenhadas como um percurso de
diagnóstico: da refinaria para a unidade, da unidade para o vetor, do
vetor para a \emph{baseline} que o julga. Todas partilham o cabeçalho
--- métrica, janela, ano e mês, ano de referência da \emph{baseline}, data
da última atualização --- e a mesma semântica: um desvio é sempre medido
contra a \emph{baseline} do ano de referência selecionado, aplicada à carga
do próprio dia, e é positivo quando o consumo excede o previsto.
%% FIGURA: prints_dashboard (uma por pagina), recortadas ao canvas.

\paragraph{Vista comparativa da refinaria.} A primeira página responde à
pergunta da reunião da manhã: onde olhar primeiro. Uma matriz de
semáforos cruza as 16 unidades com os seis vetores e mostra o desvio na
métrica e janela escolhidas, distinguindo ``sem vetor'', ``sem
\emph{baseline}'' e ``sem dados na janela''; um cartão de extremos
identifica a melhor e a pior unidade e o melhor e o pior par
unidade--vetor; um gráfico de barras ordena as unidades pelo desvio; um
diagrama de Pareto por par unidade--vetor mostra que poucos pares
concentram o excesso acumulado; um cartão de excesso acumulado dá o total
do período, o ritmo a 7 e a 30 dias contra a \emph{baseline}, e as unidades
que mais contribuem; a tendência mensal do excesso da refinaria contra a
\emph{baseline} mensal; um cartão dos piores e melhores dias do período
com o par dominante em cada um; e um cartão de qualidade das
\emph{baselines} do ano de referência --- coeficiente de determinação
médio, contagem de relações adequadas, marginais e inadequadas, e a
qualidade por unidade. Este último cartão é a primeira forma em que o
diagnóstico do Capítulo~\ref{cha:diagnostico} entrou na plataforma: ainda
antes de qualquer teste, mostra ao utilizador que a maioria das relações
com a carga explica pouco. As ordenações desta página descrevem o
comportamento atual e herdam os limites das \emph{baselines}: a ``pior
unidade'', o Pareto do excesso e os semáforos ordenam desvios face à
referência, que misturam desempenho energético com erro do modelo, escala
da unidade e número de dias, e um coeficiente de determinação médio não
certifica a qualidade de nenhuma referência. Uma unidade com uma referência
sem conteúdo preditivo pode aparecer como a pior só por isso; saber o
estado da referência muda a decisão de onde olhar primeiro. A
Secção~\ref{sec:implicacoes-plataforma} propõe mostrar o maior desvio
observado com o domínio, o volume processado e a incerteza ao lado.

\paragraph{Detalhe da unidade.} A segunda página fixa uma unidade e mostra
os seus vetores: cartões de melhor e pior vetor, índice de intensidade,
GNE total e específico do dia e acumulado do período com variação face à
\emph{baseline}, CO$_2$ equivalente, e excesso acumulado decomposto por
vetor com o seu ritmo. Um gráfico de linhas sobrepõe o consumo total ao
consumo específico com as médias móveis de 7 e 30 dias; um gráfico de
barras a 100\% mostra, por vetor, a fração de dias abaixo e acima da
banda de $\pm2\sigma$; um mapa de calor de calendário --- um quadrado por
dia, colorido pelo resíduo padronizado --- torna visíveis os períodos
persistentes e as sequências, com uma sobreposição opcional do plano de
fabrico para ligar os dias a campanhas; e uma matriz de prioridades coloca
cada vetor num plano de fração de dias em alarme contra desvio médio
fora da banda, com o tamanho proporcional ao excesso acumulado. É a
página que responde ``qual vetor, e há quanto tempo''.

\paragraph{Índice de intensidade.} A terceira página apresenta o índice
de intensidade energética da unidade (Secção~\ref{sec:metricas}) no dia e
nas janelas de 7 e 30 dias, sobre as bandas de quartil da indústria; os
cartões de índice atual e médio, do factor padrão da unidade, do quadrante
e do peso da unidade no índice da refinaria; uma cascata que decompõe o
índice por vetor e pela \emph{baseline}; e o consumo real contra a energia
padrão dia a dia, com o contexto de carga e de qualidade da
\emph{baseline} ao lado. Serve para situar a unidade face a uma referência
externa que a refinaria já usa; não substitui a \emph{baseline} interna nem
foi objeto do diagnóstico.

\paragraph{\emph{Baselines} e diagnóstico.} A quarta página é a que a
avaliação do Capítulo~\ref{cha:diagnostico} tomou como objeto. Para um
vetor, uma unidade e um ano de referência: o diagrama de dispersão do
consumo contra a carga, com a reta estimada, a banda de $\pm2\sigma$ e os
pontos coloridos pelo tempo, do mais antigo ao mais recente; um cartão de
estatísticas com a equação, o coeficiente de determinação no ano de
referência e no período, o desvio-padrão residual, a fração de dias acima
e abaixo da banda; um painel com as oito regras de carta de controlo e o
número de ocorrências de cada uma, o total de dias limpos e com violação e
a regra dominante; a soma acumulada dos resíduos desde o início do
período, com o ritmo a 7 e a 30 dias, que mostra um desvio persistente
antes de a banda diária o assinalar; e os resíduos padronizados no tempo
com as bandas de 1, 2 e 3 desvios-padrão, ou, por alternância, o consumo
observado contra o previsto. No estado de julho de 2026, para o fuel gás
da destilação atmosférica com referência em 2024, o painel mostrava sete
das oito regras violadas e 124 dias em 295 com pelo menos uma violação, a
regra dominante a ser oito dias consecutivos do mesmo lado --- o retrato,
já visível no instrumento, do que o Capítulo~\ref{cha:diagnostico} mede
com inferência adequada.

\paragraph{O que a monitorização passou a permitir.} Antes da plataforma,
o consumo por unidade e vetor existia em relatórios mensais por fábrica,
fechados algumas semanas depois do mês, sem normalização pela carga e sem
comparação entre unidades além da leitura manual. Com ela, o mesmo consumo
está disponível ao dia, para as 16 unidades e os sete vetores, com a
carga, a \emph{baseline}, o resíduo e o acumulado no mesmo lugar, e com a
proveniência de cada valor registada --- que canal o forneceu, e se um
poder calorífico ou uma propriedade foi medido ou assumido. A pergunta
``esta unidade consome mais do que devia?'' passou a ter uma resposta
consultável; a pergunta seguinte --- se essa resposta é estatisticamente
válida --- é a do capítulo seguinte.

\section{Cobertura e limitações da plataforma}
\label{sec:cobertura}

\paragraph{O que está dentro.} Dezasseis unidades processuais, sete
vetores, seis anos e sete meses de história diária, três canais
reconciliados, \emph{baselines} por ano e vetor materializadas e
exportáveis, índice de intensidade por unidade, planos de fabrico desde
2020 e propriedades da carga onde o historiador as tem.

\paragraph{O que ficou de fora, por decisão.} O gás natural e o
\emph{tail gas} não entram no perímetro de nenhuma unidade. O coque, por
ser subproduto queimado no regenerador e não combustível orçamentável,
tem \emph{baseline} própria mas fica fora do agregado global. As secções
sem casa própria no \emph{dashboard} (as secções auxiliares agregadas na
composição de cada unidade) contam para a unidade que as contém, com a
composição declarada na dimensão. Os \emph{flares} não têm
\emph{baseline}, porque não têm carga.

\paragraph{O que ficou de fora, por limitação dos dados.} A energia
elétrica é uma alocação mensal distribuída uniformemente pelos dias:
qualquer leitura diária dela é uma leitura do mês, e o
Capítulo~\ref{cha:diagnostico} trata-a como tal. Uma parte das
\emph{tags} do historiador referidas no mapeamento está anotada como ``a
verificar'' e o valor correspondente vem do balanço; onde uma
propriedade da carga não tem \emph{tag}, o modelo usa o valor por omissão
declarado e regista-o como assumido. Os poderes caloríficos sem
\emph{tag} são propagados do último valor conhecido quando um mês não os
traz, com a origem registada. A resolução é diária: transientes de
arranque, perturbações de turno e alterações rápidas de utilidades não
são visíveis --- o instrumento é tático e supervisório, não de sala de
controlo. Não há calendário operacional integrado (paragens, arranques,
intervenções), o que o plano de fabrico sobreposto ao calendário mitiga
mas não substitui; a consequência para a avaliação está na
Secção~\ref{sec:detetabilidade}.

\paragraph{O que a validação deixou aberto.} A definição de carga para os
indicadores (com ou sem resíduo reprocessado; hidrocarbonetos ao nível do
componente ou linha total do balanço) é uma decisão pendente que muda os
consumos específicos; duas fórmulas de alocação foram devolvidas ao
controlo processual; o fuel gás está validado a duas fontes dependentes; e
o mecanismo de zeros fabricados no canal de balanço, que a estimação de
\emph{baselines} exclui por filtro, continua por resolver na fonte.

\paragraph{O que o instrumento não decide.} A plataforma mostra desvios
face a uma \emph{baseline}; não diz se a \emph{baseline} é válida, nem se um
desvio é ineficiência. A primeira pergunta é a do
Capítulo~\ref{cha:diagnostico}; a segunda exige informação que a
refinaria tem e a plataforma ainda não recebe --- eventos operacionais,
rótulos, margens de materialidade --- e a Secção~\ref{sec:implicacoes-plataforma}
diz como o desenho fica preparado para a receber.
